% Abstract page (1 page). Compiled by abstract.tex (alone) and main.tex (full proposal).
% Rewritten 2026-11-24 in Dylan's style (nTof_x17/papers/WRITING_STYLE.md); beam request table added (template style).
\begin{center}
\begin{large}
{\bf AN2000 ACCELERATOR}\\
\vspace{.4cm}
{\bf Feasibility study of an X17 search in \Li\pee\Be{}\\ with the MX17 Micromegas-TPC spectrometer}\\
\vspace{.2cm}
{\bf Experimental Setup: MX17, two-arm configuration}
\end{large}
\vspace{.4cm}
\begin{flushleft}
{\bf Spokesperson(s): \todo{to decide: D.~Neff / P.~Mastinu}}\\
{\bf GLIMOS: \todo{local, to decide}}
\end{flushleft}

D.~Neff$^1$, P.~Mastinu$^2$, \todo{C.~Massimi, A.~Mengoni, G.~Gervino, C.~Gustavino, L.~Cosentino, \dots: names and affiliations to confirm}
\vspace{.4cm}

\textit{\small
$^1$ \todo{affiliation}. $^2$ INFN, Laboratori Nazionali di Legnaro, Legnaro (Padova), Italy.\\
\todo{others}
}
\end{center}
\vspace{0.65cm}

\thispagestyle{empty}
\pagestyle{empty}

\begin{center}
{\bf Abstract}
\end{center}

The ATOMKI group reported an excess of $e^+e^-$ pairs at large opening angles in the decay of the 18.15~MeV
$1^+$ state of \Be{}. They interpret the excess as a new boson with a mass of about 17~MeV, called \Xs{}.
The anomaly is still neither confirmed nor excluded. A two-arm spectrometer in Hanoi reported a similar
excess, while MEG~II found none. In addition, a 2026 study showed that cosmic rays alone can produce a bump
near $140^\circ$ in ATOMKI-type spectrometers.

The MX17 apparatus was built for, and operated in 2026 at, CERN n\_TOF. Each arm tracks the lepton in a large
Micromegas time projection chamber, placed in front of a scintillator wall, plastic scintillators and a liquid
scintillator. For this reason MX17 measures the pair opening angle from the hit positions, and it rejects
cosmic rays by their straight tracks that cross two arms. With four arms it covers about 40\,\% of the solid
angle for the \Xs{} pairs, while scintillator telescopes cover a few per cent.

We propose an 8-day feasibility study on AN2000, with two opposite arms of the apparatus around a thin
metallic \Li{} target. We will measure (i)~the rates, the live time and the detector occupancy on a proton
beam; (ii)~the acceptance, the angular resolution and the energy scale on the 17.64~MeV M1 transition at the
441~keV resonance, where no anomaly is expected; (iii)~the cosmic-ray and beam-induced backgrounds, with and
without the tracking veto; and (iv)~the 18.15~MeV region at $E_p = 1.10$~MeV. For item (iv) we will use the
existing copper target holder. With this holder, a Geant4 simulation gives an expected 90\,\% CL limit in five
days of $2.35\times R_\mathrm{ATOMKI}$, which is close to the MEG~II limit ($2.07\times R_\mathrm{ATOMKI}$). A
lighter aluminium end piece would bring the expected limit to about $0.6$--$0.8\times R_\mathrm{ATOMKI}$, below the
ATOMKI ratio (an estimate scaled from the four-arm simulation). The results
will fix the design and the beam request of the full four-arm experiment.

\vspace{.5cm}
\begin{center}
{\bf Beam request}\\[.2cm]
\begin{tabular}{|c|c|c|c|c|c|c|}
\hline
\ Beam & $E_p$ (MeV) & Intensity & Beam mode & Target & Backing & Beam time \ \\
\hline
\ $p$ & 0.441 & $\le 1$~\si{\micro\ampere} & continuous & metallic \Li{}, 10--12~\si{\micro\metre} & Cu, 300~\si{\micro\metre} & 1 day \ \\
\ $p$ & 1.10 & $\le 1$~\si{\micro\ampere} & continuous & metallic \Li{}, 10--12~\si{\micro\metre} & Cu, 300~\si{\micro\metre} & 7 days \ \\
\hline
\multicolumn{6}{|r|}{total \ } & 8 days \ \\
\hline
\end{tabular}\\[.15cm]
Experimental setup: MX17, two arms. Accelerator: AN2000, in the second half of 2027.
\end{center}
